\section{Conclusão}\label{sec:conclusion}

Este artigo quantificou o ganho de PTF que neutralizaria a perda de produto de curto prazo associada a tetos alternativos de jornada. A curva simulada é acentuadamente não linear: partindo da distribuição habitual formal limitada a 44h e sob a calibração central de eficiência, o teto de 40 horas requer $1{,}53$--$1{,}57\%$ de PTF, enquanto a passagem para 36 horas requer entre $6{,}52\%$ e $8{,}38\%$. A redução física das horas formais é a maior parcela negativa da decomposição; a resposta da eficiência depende da especificação, enquanto a realocação para o emprego informal amplia a perda.

Os contrafactuais também separam três objetos que não devem ser confundidos: a perda de produto com PTF fixa, o ganho de PTF que recompõe essa perda e o indicador GHH de consumo e lazer. Reduções moderadas podem diminuir o produto e, simultaneamente, elevar o indicador representativo de bem-estar. No teto de 36 horas, o requisito de PTF é quantitativamente maior e o equivalente de consumo é negativo nas duas funções de eficiência, variando de $-3{,}42\%$ a $-1{,}66\%$. As sensibilidades mostram que a avaliação depende da fadiga e da distribuição inicial de horas. A base de 44h é uma hipótese de mensuração da reforma, distinta de horas contratadas observadas.

As conclusões são condicionais a capital e emprego agregados predeterminados. Incorporar investimento, barganha, seleção de trabalhadores, ligações setoriais e respostas de oferta de trabalho é a extensão natural para uma avaliação completa. Dentro do objeto estudado, os resultados mostram que a intensidade da redução e a resposta da produtividade alteram de forma substantiva sua aritmética macroeconômica e motivam a avaliação de trajetórias graduais e de instrumentos de transição.

\FloatBarrier
